Il monitoraggio continuo dell'elettrocardiogramma impone di classificare ogni battito mentre il segnale arriva, rispettando vincoli di latenza e affidabilità.
Molti lavori sulla classificazione automatica si limitano all'addestramento e alla valutazione offline, senza collegare ingestion, inferenza e persistenza in un'architettura eseguibile.

Si propone una pipeline end-to-end per la classificazione \emph{near-real-time} di battiti ECG, con approccio ibrido batch/streaming.
Un \texttt{RandomForestClassifier}, addestrato in batch su feature a livello di battito dal dataset MIT-BIH in formato tabulare, viene distribuito come artefatto e applicato in streaming da Apache Spark Structured Streaming su eventi Kafka pubblicati da un simulatore edge.
Le predizioni sono scritte su Delta Lake con checkpointing, garantendo recupero dopo failure e disaccoppiamento tra producer e consumer.
L'implementazione in Scala è riproducibile con Docker e \texttt{sbt}.

Sul test set si osservano accuratezza globale circa al 96\% e latenza end-to-end tipicamente sub-secondo, con limiti sulle classi rare e implicazioni cliniche degli errori di classificazione analizzate nel contesto delle super-classi AAMI.
Il significato del lavoro è dimostrare che Kafka, Spark e Delta Lake possono comporsi in un sistema modulare per dati biomedici strutturati, fungendo da template per pipeline e-Health che uniscano qualità predittiva e proprietà non funzionali di un deployment distribuito.
